\begin{musica}{Olá Tuna}{António Mafra}{Tuna de Engenharia da Universidade do Porto}
\nota{Introdução: C G C}

De \ac{C}{o}lhos brilhantes
Somos os tunantes
Com tudo p’ra ser fe\ac{G7}{li}zes.
Porém com chouriço
Cerveja é preciso,
Disso somos bons a\ac{C}{man}tes.

As cordas chorando,
No peito batendo
Paixão num beijo sedento.
Cantar a valer, pena de saber
O fim de um eterno momento.

\nota{Refrão:}
\ac{C}{Quan}do passa, quando passa
Ninguém dela os olhos \ac{G7}{ti}ra,
Ai que graça, ai que graça
A Tuna de Engenha\ac{C}{ria}.

Cantando baladas,
Fazendo noitadas,
Serenatas às donzelas.
De rosa na mão,
Dar o coração
Fazendo tudo por elas.

Se ela sorriu
Gostou do que ouviu
Mas nos achou desordeiros.
Achamos banal,
Mesmo natural,
Pois nós somos engenheiros.

\nota{Refrão}

\posicoes{C: x32010, G7: 320001}
\end{musica}
